\documentclass[11pt]{article}
\usepackage[margin=1in]{geometry}
\usepackage[T1]{fontenc}
\usepackage{lmodern}
\usepackage{amsmath,amssymb,booktabs,longtable,hyperref}
\hypersetup{colorlinks=true,linkcolor=blue,urlcolor=blue}

\title{Dimension-5 ORBIT Comparison to a Sign-Flipped Einstein-Hilbert Reference}
\author{ORBIT}
\date{2026-03-17}

\begin{document}
\maketitle
\tableofcontents

\section{Setup}
This report records the final lightweight comparison run carried out after introducing an explicit user-controlled kinetic-sign convention into the ORBIT comparison layer. The comparison was intentionally restricted to mass dimension
\[
d \le 5
\]
in order to avoid the substantially more expensive full \(d \le 6\) rerun.

The conventions used for this final run are:
\begin{itemize}
\item mostly-minus Lorentzian signature;
\item metric perturbation \(g_{\mu\nu} = \eta_{\mu\nu} + \kappa h_{\mu\nu}\);
\item reference action
\[
\mathcal{L}_{\mathrm{ref}}
=
\Lambda_{\mathrm{ref}} \sqrt{-g}
- \frac{2}{\kappa_{\mathrm{ref}}^2}\sqrt{-g}\,R;
\]
\item canonical normalization of the Matchete input to the predetermined quadratic target sign
\[
s_{\mathrm{kin}}=-1.
\]
\end{itemize}

The Matchete input is normalized before comparison by setting \(hbar \to 1\) and extracting the finite \(\epsilon^0\) part term by term.

\section{Canonical Normalization}
The quadratic derivative sector of the Matchete result is proportional to the Fierz--Pauli kinetic term,
\[
\mathcal{L}^{\mathrm{Matchete}}_{2,2}
=
Z_h\,\mathcal{L}_{\mathrm{FP}},
\qquad
Z_h=\frac{M^2\kappa^2\bigl(1+\log(\mu^2/M^2)\bigr)}{24}.
\]
For the present run, the canonicalization target is chosen directly by user convention rather than inferred from the reference branch. The field rescaling therefore uses
\[
\beta^2=\frac{s_{\mathrm{kin}}}{Z_h}=-\frac{1}{Z_h},
\]
so that the quadratic sector is mapped to \(-\mathcal{L}_{\mathrm{FP}}\).

\section{Sector Support}
At \(d \le 5\), the supported reference sectors are
\[
\{1,0\},\ \{2,0\},\ \{2,2\},\ \{3,0\},\ \{3,2\},\ \{4,0\},\ \{5,0\}.
\]
The final comparison found no unsupported leftover sectors:
\[
\texttt{UnsupportedSectors} = \varnothing,
\qquad
\texttt{UnsupportedDifference}=0.
\]

\section{Reduction Summary}
\subsection{Reduced Matchete Input}
\begin{longtable}{rrrrr}
\toprule
Sector & Input terms & IBP terms & Reduced terms & Physical coordinates \\
\midrule
\{1,0\} & 2 & 2 & 2 & 1 \\
\{2,0\} & 4 & 4 & 4 & 2 \\
\{2,2\} & 10 & 4 & 4 & 4 \\
\{3,0\} & 6 & 6 & 6 & 3 \\
\{3,2\} & 42 & 13 & 6 & 10 \\
\{4,0\} & 10 & 10 & 10 & 5 \\
\{5,0\} & 14 & 14 & 14 & 7 \\
\bottomrule
\end{longtable}

\subsection{Reduced Reference}
\begin{longtable}{rrrrr}
\toprule
Sector & Input terms & IBP terms & Reduced terms & Physical coordinates \\
\midrule
\{1,0\} & 1 & 1 & 1 & 1 \\
\{1,2\} & 2 & 0 & 0 & 0 \\
\{2,0\} & 2 & 2 & 2 & 2 \\
\{2,2\} & 10 & 4 & 4 & 4 \\
\{3,0\} & 3 & 3 & 3 & 3 \\
\{3,2\} & 28 & 13 & 6 & 10 \\
\{4,0\} & 5 & 5 & 5 & 5 \\
\{5,0\} & 6 & 6 & 6 & 7 \\
\bottomrule
\end{longtable}

\section{Solved Matching Constraints}
The final saved result satisfies
\[
\texttt{ExactMatchQ} = \mathrm{True},
\qquad
\texttt{HasSolution} = \mathrm{True}.
\]

The explicit solved branches in the stored result reduce to the following invariant conditions:
\begin{align}
\mu^2 &> 0, \\
M^2 &= e^{3/2}\mu^2, \\
\Lambda_{\mathrm{ref}} &= 0, \\
\kappa_{\mathrm{ref}}^2 &= \frac{48}{M^2}.
\end{align}
In addition, the sign of \(\kappa_{\mathrm{ref}}\) is correlated with the sign of the input \(\kappa\):
\[
\kappa>0 \Rightarrow \kappa_{\mathrm{ref}}>0,
\qquad
\kappa<0 \Rightarrow \kappa_{\mathrm{ref}}<0.
\]

Equivalently, the saved \texttt{Reduce} output is
\begin{verbatim}
\[Mu]bar2 > 0 &&
((M == E^(3/4)*Sqrt[\[Mu]bar2] &&
  ((\[Kappa] > 0 && LambdaRefFinal == 0 &&
    KappaRefFinal == 4*Sqrt[3]*Sqrt[M^(-2)]) ||
   (\[Kappa] < 0 && LambdaRefFinal == 0 &&
    KappaRefFinal == -4*Sqrt[3]*Sqrt[M^(-2)]))) ||
 (M == -(E^(3/4)*Sqrt[\[Mu]bar2]) &&
  ((\[Kappa] > 0 && LambdaRefFinal == 0 &&
    KappaRefFinal == 4*Sqrt[3]*Sqrt[M^(-2)]) ||
   (\[Kappa] < 0 && LambdaRefFinal == 0 &&
    KappaRefFinal == -4*Sqrt[3]*Sqrt[M^(-2)]))))
\end{verbatim}

\section{Conclusion}
After fixing the canonical-sign logic so that the kinetic target sign is an explicit user-controlled convention, the ORBIT comparison closes cleanly at \(d \le 5\) in the sign-flipped Einstein--Hilbert branch. In this final run:
\begin{itemize}
\item the comparison is exact in the supported sector space;
\item no unsupported sectors survive;
\item the cosmological constant is forced to vanish;
\item the remaining parameter relations are the ones listed above.
\end{itemize}

The expensive \(d \le 6\) rerun was intentionally deferred. The current evidence therefore supports the statement:
\begin{quote}
the canonically normalized Matchete result matches a sign-flipped Einstein--Hilbert reference through mass dimension five, once the user-chosen kinetic sign is fixed to \(s_{\mathrm{kin}}=-1\).
\end{quote}

\section{Primary Artifact}
The machine-readable source for this report is the saved comparison result
\begin{verbatim}
orbit_fixed_eh_comparison_d5_final_result.wl
\end{verbatim}
which contains the full reduced input, reduced reference, solved constraints, and sector-by-sector comparison data.

\end{document}
